% 原书物理页=印刷页4--7；1.2 全节.
\section{经典群}\label{sec:lie-1-2}
作为定理~\ref{thm:lie-1-1-8} 的另一个有趣的推论，所有``\newterm{lie-classical-group}{经典群}{classical group}''都是 Lie 群. 由于它们是最有意思的例子，我们在下面列出其清单. 字母 $K$ 统一指代域 $\mathbb R$ 与 $\mathbb C$，而 $V$ 指一个维数为有限数 $n$ 的 $K$-向量空间.
\begin{numpara}[特殊线性群]\label{np:lie-1-2-1}
行列式映射 $\mathrm{GL}(V)\to\mathrm{GL}_1(K)$ 是拓扑群之间的连续同态，其核 $\mathrm{SL}(V)$ 是闭的，因而是 Lie 子群. 我们称之为 $V$ 的\newterm{lie-special-linear-group}{特殊线性群}{special linear group}.$V$ 的一组基的选取将其等同于 $\mathrm{SL}_n(K)$.
\end{numpara}
\begin{exercise*}
$\mathrm{SL}_n(K)$ 的中心由行列式为 $1$ 的\newterm{lie-homothety}{位似}{homothety}组成，因此它同构于 $K^*$ 中 $n$ 次单位根所成的群 $\mu_n(K)$（当 $K=\mathbb C$ 时其阶为 $n$；当 $K=\mathbb R$ 时阶为 $1$ 或 $2$）.
\end{exercise*}
\begin{numpara}[正交群]\label{np:lie-1-2-2}
记 $\mathrm{Sym}_K(V)$ 为 $K$-\newterm{lie-symmetric-bilinear-form}{对称双线性型}{symmetric bilinear form} $\Phi:V\times V\to K$ 构成的 $K$-向量空间.$\mathrm{GL}(V)$ 在 $\mathrm{Sym}_K(V)$ 上有一个线性作用（从而是连续作用），由 $g\Phi(v,w):=\Phi(g^{-1}v,g^{-1}w)$ 给出. 于是 $\Phi$ 在该作用下的\newterm{lie-stabilizer}{稳定子群}{stabilizer}在 $\mathrm{GL}(V)$ 中是闭的，因而是 Lie 子群. 我们把它记作
\[
\mathrm{O}(\Phi):=\{g\in\mathrm{GL}(V),\ \Phi(gv,gw)=\Phi(v,w),\ \forall v,w\in V\}.
\]
若 $\Phi$ 非退化，则称之为 $\Phi$ 的\newterm{lie-orthogonal-group}{正交群}{orthogonal group}.
回顾一下，两个\newterm{lie-quadratic-form}{二次型}{quadratic form}称为\emph{等价}的，如果它们属于 $\mathrm{GL}(V)$ 的同一轨道. 由此可知，它们的正交群在 $\mathrm{GL}(V)$ 中共轭，特别地彼此同构. 当 $K=\mathbb C$ 时，我们知道所有二次型都等价，从而 $\mathrm{O}(\Phi)$ 同构于通常的复正交群
\[
\mathrm{O}(n,\mathbb C)=\{M\in\mathrm{M}_n(\mathbb C),\ ^tMM=I_n\}.
\]
当 $K=\mathbb R$ 时，Sylvester 定理断言：存在整数 $p\leqslant n$ 与一组基 $e_1,\cdots,e_n$，使得
\[
\Phi\Bigl(\sum_{i=1}^{n}x_ie_i,\ \sum_{i=1}^{n}y_ie_i\Bigr)=\sum_{i=1}^{p}x_iy_i-\sum_{i=p+1}^{n}x_iy_i.
\]
整数对 $(p,\ q:=n-p)$ 称为 $\Phi$ 的\newterm{lie-signature}{符号差}{signature}.于是 $\mathrm{O}(\Phi)$ 同构于矩阵群
\[
\mathrm{O}(p,q):=\{M\in\mathrm{M}_{p+q}(\mathbb R),\ ^tM.D_{p,q}.M=D_{p,q}\}
\]
其中 $D_{p,q}=D(1,\cdots,1,-1,\cdots,-1)$ 是对角矩阵，$1$ 重复 $p$ 次，$-1$ 重复 $q$ 次. 当 $p=n$ 与 $p=0$ 时，我们重新得到欧氏正交群 $\mathrm{O}(n,0)=\mathrm{O}(n)$，众所周知它是紧的. 其余的正交群都不紧. 我们有 $\mathrm{O}(p,q)=\mathrm{O}(q,p)$，而且可以证明，当 $q\leqslant p$ 时诸群 $\mathrm{O}(p,q)$ 两两不同构.
\end{numpara}
\begin{example*}
对时空 $\mathbb R^4$ 上的型 $xx'+yy'+zz'-tt'$，我们由此得到广义相对论中出现的 \newterm{lie-lorentz-group}{Lorentz 群}{Lorentz group} $\mathrm{O}(3,1)$.
\end{example*}
此外，我们定义特殊正交群为
\[
\mathrm{SO}(p,q):=\mathrm{O}(p,q)\cap\mathrm{SL}_{p+q}(\mathbb R).
\]
\begin{numpara}[酉群]\label{np:lie-1-2-3}
这里 $K=\mathbb C$. 记 $\mathrm{Herm}(V)$ 为\newterm{lie-hermitian-form}{埃尔米特型}{hermitian form} $\Psi:V\times V\to\mathbb C$ 构成的 $\mathbb R$-向量空间.$\mathrm{GL}(V)$ 在 $\mathrm{Herm}(V)$ 上有一个线性作用（从而是连续作用），由 $g\Psi(v,w):=\Psi(g^{-1}v,g^{-1}w)$ 给出. 于是 $\Psi$ 在该作用下的稳定子群在 $\mathrm{GL}(V)$ 中是闭的，因而是 Lie 子群. 我们把它记作
\[
\mathrm{U}(\Psi):=\{g\in\mathrm{GL}(V/\mathbb C),\ \Psi(gv,gw)=\Psi(v,w),\ \forall v,w\in V\}.
\]
若 $\Psi$ 非退化，则称之为 $\Psi$ 的\newterm{lie-unitary-group}{酉群}{unitary group}\footnote{译注：$\mathrm{GL}(V/\mathbb C)$ 为原文记号，指把 $V$ 视为 $\mathbb C$-向量空间的一般线性群.}.
与对称情形一样，两个埃尔米特型称为等价的，如果它们属于 $\mathrm{GL}(V)$ 的同一轨道. 此时，它们的酉群同构，甚至共轭. 这里同样，埃尔米特型由其符号差分类：Sylvester 定理断言，存在整数 $p\leqslant n$ 与一组基 $e_1,\cdots,e_n$，使得
\[
\Psi\Bigl(\sum_{i=1}^{n}x_ie_i,\ \sum_{i=1}^{n}y_ie_i\Bigr)=\sum_{i=1}^{p}x_i\overline{y}_i-\sum_{i=p+1}^{n}x_i\overline{y}_i.
\]
于是，令 $q=n-p$，拓扑群 $\mathrm{U}(\Psi)$ 同构于矩阵群
\[
\mathrm{U}(p,q):=\{M\in\mathrm{M}_{p+q}(\mathbb C),\ ^t\overline{M}.D_{p,q}.M=D_{p,q}\}
\]
当 $pq=0$ 时，我们重新得到通常的酉群 $\mathrm{U}(n,0)=\mathrm{U}(n)$，如我们所知它是紧的. 其余的酉群都不紧. 我们有 $\mathrm{U}(p,q)=\mathrm{U}(q,p)$，而且可以证明，当 $q\leqslant p$ 时诸群 $\mathrm{U}(p,q)$ 两两不同构.
\end{numpara}
\begin{remark*}
群 $\mathrm{U}(1,n-1)$ 在当今的数论问题中起着主导作用.
\end{remark*}
与前面一样，我们有特殊酉群
\[
\mathrm{SU}(p,q):=\mathrm{U}(p,q)\cap\mathrm{SL}_{p+q}(\mathbb C).
\]
\begin{exercise*}
确定 $\mathrm{U}(n)$、$\mathrm{SU}(n)$、$\mathrm{O}(n)$、$\mathrm{SO}(n)$ 的中心.
\end{exercise*}
\begin{numpara}[辛群]\label{np:lie-1-2-4}
这里 $K$ 仍表示 $\mathbb R$ 或 $\mathbb C$. 记 $\mathrm{Alt}_K(V)$ 为 $K$-\newterm{lie-alternating-form}{交错双线性型}{alternating bilinear form} $\psi:V\times V\to K$ 构成的 $K$-向量空间.$\mathrm{GL}(V)$ 在 $\mathrm{Alt}_K(V)$ 上有一个线性作用（从而是连续作用），由 $g\psi(v,w):=\psi(g^{-1}v,g^{-1}w)$ 给出. 于是 $\psi$ 在该作用下的稳定子群在 $\mathrm{GL}(V)$ 中是闭的，因而是 Lie 子群. 我们把它记作
\[
\mathrm{Sp}(\psi):=\{g\in\mathrm{GL}(V),\ \psi(gv,gw)=\psi(v,w),\ \forall v,w\in V\}.
\]
若 $\psi$ 非退化（此时称 $\psi$ 为\newterm{lie-symplectic-form}{辛形式}{symplectic form}），则称之为 $\psi$ 的\newterm{lie-symplectic-group}{辛群}{symplectic group}.此时 $V$ 的维数必为偶数，我们改用记号 $\dim_K(V)=2n$. 同前面一样，这个群在同构的意义下只依赖于 $\psi$ 的等价类. 然而我们知道，所有辛形式都是等价的. 特别地，存在 $V$ 的一组基 $e_1,\cdots,e_{2n}$，使得
\[
\psi\Bigl(\sum_{i=1}^{n}x_ie_i,\ \sum_{i=1}^{n}y_ie_i\Bigr)=\sum_{i=1}^{n}(x_iy_{n+i}-x_{n+i}y_i).
\]
换言之，$\psi$ 在这组基下的矩阵是反对称矩阵 $J=J_{2n}=\begin{pmatrix}0_n & I_n\\ -I_n & 0_n\end{pmatrix}$. 于是我们看到，$\mathrm{Sp}(\psi)$ 同构于辛群
\[
\mathrm{Sp}_{2n}(K)=\{M\in\mathrm{GL}_{2n}(K),\ ^tMJM=J\}.
\]
于是 $\mathrm{Sp}_2(K)=\mathrm{SL}_2(K)$. 更一般地，可以证明辛矩阵的行列式总等于 $1$，从而对一切 $n$，$\mathrm{Sp}_{2n}(K)$ 是 $\mathrm{SL}_{2n}(K)$ 的闭的非紧子群.
\end{numpara}
\begin{exercise*}
证明：$\mathbb C$-向量空间 $V$ 上一个埃尔米特型的实部（相应地，虚部）是其基础 $\mathbb R$-向量空间 $\widetilde{V}$ 上的 $\mathbb R$-双线性对称型（相应地，交错型）. 由此推得：在 $\mathrm{GL}_n(\mathbb C)\simeq\{P\in\mathrm{GL}_{2n}(\mathbb R),\ P^{-1}J_{2n}P=J_{2n}\}\subset\mathrm{GL}_{2n}(\mathbb R)$ 中，有 $\mathrm{U}(n)=\mathrm{O}(2n)\cap\mathrm{Sp}_{2n}(\mathbb R)$.
\end{exercise*}
\begin{numpara}[四元数群]\label{np:lie-1-2-5}
回顾一下，Hamilton 的\newterm{lie-quaternion}{四元数}{quaternion}体 $\mathbb H$ 是中心为 $\mathbb R$ 的\newterm{lie-division-algebra}{可除代数}{division algebra}，它有一组 $\mathbb R$-基 $1,i,j,k$，乘法在其中由 $ij=k$、$jk=i$、$ki=j$、$i^2=j^2=k^2=1$\footnote{译注：按四元数的标准定义应有 $i^2=j^2=k^2=-1$，原文如此.} 给出. 通过 $i\mapsto i$ 将 $\mathbb C$ 送入 $\mathbb H$，也可以把 $\mathbb H=\mathbb C\oplus\mathbb Cj\simeq\mathbb C^2$ 看成一个 $2$ 维 $\mathbb C$-向量空间，但其乘法不是 $\mathbb C$-线性的.
把 $\mathbb H^n$ 视为 $\mathbb H$ 上的右向量空间，并令 $\mathrm{GL}_n(\mathbb H)$ 为 $\mathbb H^n$ 的 $\mathbb H$-线性自同构组成的群. 借助等同 $\mathbb H^n=\mathbb C^n+j\mathbb C^n$，有 $\mathrm{GL}_n(\mathbb H)=\{P\in\mathrm{GL}_{2n}(\mathbb C),\ PJ_{2n}=J_{2n}\overline{P}\}$. 因此它是 $\mathrm{GL}_{2n}(\mathbb C)$ 的闭子群，从而是 Lie 群.
也可以在 $\mathbb H$ 上装备\newterm{lie-anti-involution}{反对合}{anti-involution} $\sigma$，它由 $\sigma(i)=-i$、$\sigma(j)=-j$、$\sigma(k)=-k$ 给出，并在 $\mathbb H^n$ 上装备典范的 $\sigma$-埃尔米特型. 这一型的各向同性子群 $\mathrm{U}_n(\mathbb H)\subset\mathrm{GL}_n(\mathbb H)$ 仍是 Lie 群，并且如上面习题所示，可以把它等同于 $\mathrm{U}(2n)\cap\mathrm{Sp}_{2n}(\mathbb C)$.
\end{numpara}
\begin{numpara}[经典群的重要性]\label{np:lie-1-2-6}
它们出现在相当多样的领域之中.
\textit{几何.}\ 这是它们最钟爱的领域. 这些群常常是一些特殊空间（球面、双曲面……）的对称群，因而作为非欧几何的模型.
\textit{数论.}\ 这是一个稍显出人意料的应用，但近百年来，人们一直在形如 $G(\mathbb Z)\backslash G(\mathbb R)/K$ 的空间上研究函数空间、微分形式等，其中 $G$ 是一个实经典群，$K$ 是 $G$ 的一个\newterm{lie-maximal-compact-subgroup}{极大紧子群}{maximal compact subgroup}.例如取 $G=\mathrm{GL}_1$，便得到三角函数；取 $G=\mathrm{GL}_2$，便得到模形式；对一般的 $G$，人们称之为``自守形式''.例如，这些对象在 Fermat 定理的证明中起着关键作用.
\textit{物理.}
\begin{itemize}
\item 广义相对论的 Lorentz 群是 $\mathrm{O}(1,3)$. 它是``与时空几何相伴的''群.
\item 某些基本粒子按照 $\mathrm{SU}(3)$ 的``不可约''表示来分类.
\item 氢原子的能级由球面 $S^2$ 上球面 Laplace 算子的特征值给出；稍后我们将利用 $\mathrm{SO}_3(\mathbb R)$ 作用下 $C^\infty(S^2)$ 的``谱分解''来计算这些特征值.
\end{itemize}
\end{numpara}
